% Ficha de estudio 6 del Taller virtual ATA EPE: ruedas de fricción, bandas elásticas y manivela.
\documentclass[11pt,a4paper]{article}
% Preámbulo común de las fichas de estudio del Taller virtual ATA EPE.
\usepackage[spanish,es-noshorthands,es-tabla]{babel}
\usepackage[utf8]{inputenc}
\usepackage[T1]{fontenc}
\usepackage{lmodern}
\usepackage[margin=2.2cm]{geometry}
\usepackage[table]{xcolor}
\usepackage{amsmath,amssymb,bm}
\usepackage{tabularx,array,booktabs}
\usepackage{enumitem}
\usepackage{fontawesome5}
\usepackage[most]{tcolorbox}
\usepackage[hidelinks]{hyperref}
\usepackage{listings}

\definecolor{morado}{HTML}{2A1B7E}
\definecolor{amarillo}{HTML}{FFCF2E}
\definecolor{rojo}{HTML}{C4142C}
\definecolor{verde}{HTML}{1F9D55}
\definecolor{gris}{HTML}{5D5A73}

\newcolumntype{Y}{>{\raggedright\arraybackslash}X}
\renewcommand{\arraystretch}{1.25}
\setlist{itemsep=1pt,topsep=3pt}
\numberwithin{equation}{section}
\lstset{basicstyle=\ttfamily\small,columns=fullflexible,keepspaces=true,frame=single,rulecolor=\color{gris!40},
  literate={á}{{\'a}}1 {é}{{\'e}}1 {í}{{\'i}}1 {ó}{{\'o}}1 {ú}{{\'u}}1 {ñ}{{\~n}}1 {Ω}{{$\Omega$}}1
    {θ}{{$\theta$}}1 {λ}{{$\lambda$}}1 {Δ}{{$\Delta$}}1 {μ}{{$\mu$}}1 {·}{{$\cdot$}}1 {←}{{$\leftarrow$}}1 {§}{{\S}}1 {→}{{$\to$}}1 {≠}{{$\neq$}}1 {É}{{\'E}}1 {Á}{{\'A}}1 {∨}{{$\vee$}}1 {∧}{{$\wedge$}}1 {σ}{{$\sigma$}}1 {×}{{$\times$}}1 {≈}{{$\approx$}}1}

\newcommand{\rev}[1]{\widetilde{#1}}
\newcommand{\I}{\mathbf{I}}
\newcommand{\vect}[1]{\bm{#1}}
\newcommand{\A}{\mathcal{A}}

\newtcolorbox{meta}{colback=morado,colframe=morado,coltext=white,arc=3mm}
\newtcolorbox{lectura}{colback=morado!5,colframe=morado,fonttitle=\bfseries,title={\faBook~Lecturas}}
\newtcolorbox{ejercicio}[1]{colback=amarillo!10,colframe=amarillo!80!black,fonttitle=\bfseries,title={\faPencil*~Ejercicio #1}}
\newtcolorbox{control}{colback=rojo!5,colframe=rojo,fonttitle=\bfseries,title={\faQuestionCircle~Preguntas de control}}
\newtcolorbox{codigo}[1][De la ecuación al código]{colback=verde!6,colframe=verde,fonttitle=\bfseries,title={\faCode~#1}}
\newtcolorbox{cierre}{colback=verde!6,colframe=verde,fonttitle=\bfseries,title={\faFlagCheckered~Cierre del módulo}}
\newtcolorbox{nota}[1]{colback=amarillo!12,colframe=amarillo!80!black,fonttitle=\bfseries,title={\faLightbulb~#1}}

% \cabecera{número}{título}{subtítulo}
\newcommand{\cabecera}[3]{%
  \title{\textcolor{morado}{\textbf{Ficha #1. #2}}\\[3pt]\large Taller virtual ATA EPE · #3}
  \author{}\date{}}

\cabecera{6}{Ruedas, bandas y manivela}{Módulo 6: transmisiones}

\begin{document}
\maketitle
\vspace{-2.5em}

\begin{meta}
\textbf{\textcolor{amarillo}{Meta del módulo.}} Transmitir movimiento como lo hace el kit: ruedas que se empujan por
fricción (de canto o de forma transversal), bandas elásticas que se estiran y patinan, y una manivela que se gira con
el dedo. Cara pública: \textbf{transmitir} — el sistema de dos ruedas, bandas entre poleas, la rueda que se corre
sobre un disco y cambia de velocidad.
\end{meta}

\begin{lectura}
\begin{enumerate}[leftmargin=*]
  \item \textbf{Johnson, \emph{Contact Mechanics}, cap.~8}: rodadura, deslizamiento parcial (\emph{creep}) y
        pivoteo (\emph{spin}).
  \item \textbf{Treloar, \emph{The Physics of Rubber Elasticity}, caps.~1--4}: por qué el caucho es elástico
        (entropía) y la ley neo-Hookeana.
  \item Budynas y Nisbett, \emph{Shigley's Mechanical Engineering Design}, cap.~17: transmisiones por banda,
        tensiones y deslizamiento.
  \item Cualquier texto de estática (por ejemplo Meriam y Kraige, \emph{Statics}, fricción en bandas): la ecuación
        del cabrestante.
  \item Plan del taller, secciones de ruedas de fricción y bandas elásticas (los modelos completos).
\end{enumerate}
\end{lectura}

% =================================================================================================
\section{Rodar sin deslizar}

En el punto de contacto $P$ de dos cuerpos con velocidades $\Omega_A$ y $\Omega_B$ (en el espacio), la velocidad
relativa de las superficies sale de la ficha 1:
$\dot P_A - \dot P_B = P\times(\Omega_A - \Omega_B)$ (con el signo de la convención). Rodar es que su parte tangencial
sea nula. Con ejes fijos es una ligadura holónoma; se agrega al solucionador como una fila más.

\subsection{Canto con canto}
Ejes paralelos, radios $R_A$, $R_B$: $\omega_AR_A = -\omega_BR_B$. El sistema de dos ruedas de la cartilla
($R_A/R_B = 1/3$) gira tres veces más lento en la grande, y al revés.

\subsection{Canto contra cara (transversal)}
Una rueda de radio $R$ apoyada a distancia $r$ del eje de un disco, con ejes perpendiculares:
\begin{equation}
  \omega_{\text{rueda}}\,R = \omega_{\text{disco}}\,r \quad\Longrightarrow\quad \frac{\omega_{\text{rueda}}}{\omega_{\text{disco}}} = \frac rR .
\end{equation}
La relación cambia con la posición del contacto: es una transmisión continua.

\begin{ejercicio}{6.1}
Deduce (1.1) con la velocidad de los puntos en PGA. ¿Qué pasa con el sentido de giro de la rueda cuando el contacto
cruza el centro del disco?
\end{ejercicio}

\subsection{Deslizamiento y pivoteo}
El contacto transmite como máximo $\mu N$ de fuerza tangencial; más allá, patina. Para no tener un salto en
$\operatorname{sgn}(s)$ se regulariza: $F_t = -\mu N\tanh(s/s_0)$. En la transversal, el contacto es una franja de ancho
$w$ y sus extremos no pueden rodar a la vez: se disipa $P\approx\mu N\,\omega_{\text{disco}}\,w/4$.

\begin{ejercicio}{6.2}
Deduce la potencia de pivoteo integrando $\mu\,(N/w)\,\lvert\omega x\rvert$ para $x\in[-w/2, w/2]$.
\end{ejercicio}

% =================================================================================================
\section{El caucho}

\subsection{Por qué es elástico}
El caucho es una red de cadenas largas. Estirarlo no aumenta tanto su energía como disminuye su \emph{entropía}: las
cadenas pasan de enredadas a alineadas. La teoría de redes gaussianas da la energía de deformación
$W = \tfrac G2(\lambda_1^2+\lambda_2^2+\lambda_3^2-3)$, con $G = nk_BT$ ($n$ cadenas por volumen). Para una banda
estirada a lo largo (volumen constante, $\lambda_2 = \lambda_3 = \lambda^{-1/2}$):
\begin{equation}
  F = G\,A_0\,\bigl(\lambda - \lambda^{-2}\bigr),\qquad \lambda = L/L_0 .
\end{equation}
Para estiramientos pequeños, $\lambda = 1+\varepsilon$, se recupera Hooke con $E = 3G$.

\begin{ejercicio}{6.3}
Deduce (2.1) a partir de $W$. Verifica el límite de Hooke. ¿Por qué el caucho estirado se calienta?
\end{ejercicio}

% =================================================================================================
\section{Bandas sobre poleas}

\subsection{Geometría}
El camino de la banda está hecho de tramos rectos tangentes a dos poleas y de arcos abrazados. Entre dos círculos hay
cuatro tangentes: dos exteriores (banda abierta) y dos interiores (cruzada); cuál se usa depende del lado por el que
la banda abraza cada polea. Para dos poleas, $L$ tiene forma cerrada (plan, ec. de la banda abierta y cruzada).

\begin{ejercicio}{6.4}
Calcula las tangentes exteriores e interiores entre dos círculos del mismo plano. Programa la longitud de la banda
para $n$ poleas y verifica que con dos coincide con la fórmula cerrada.
\end{ejercicio}

\subsection{El cabrestante}
En un trozo de banda abrazado $d\theta$, la normal es $T\,d\theta$ y la fricción puede equilibrar hasta
$dT = \mu T\,d\theta$. Integrando:
\begin{equation}
  \frac{T_1}{T_2}\le e^{\mu\theta}\qquad\Longrightarrow\qquad \tau_{\max} = 2T_0R\,\frac{e^{\mu\theta}-1}{e^{\mu\theta}+1}.
\end{equation}

\begin{ejercicio}{6.5}
Deduce el par máximo de (3.1) con $T_1 = T_0 + \tau/2R$ y $T_2 = T_0 - \tau/2R$. ¿Cuánto par más transmite una banda
si se aleja una polea de modo que $T_0$ se duplique? ¿Y si se aumenta el ángulo abrazado con una polea tensora?
\end{ejercicio}

\subsection{Modelo dinámico}
La banda es un acople elástico entre los ángulos de las poleas, con deslizamiento acumulado $\delta$ cuando se supera
el límite del cabrestante (plan, modelo dinámico de la banda). Cada polea recibe el par $\mp R(T_1 - T_2)$ y su eje la
fuerza $T_1 + T_2$. Rigidez alta y masas pequeñas: hay que mirar la estabilidad ($\omega_n h < 2$, ficha 2).

% =================================================================================================
\section{La manivela y la mano}

\begin{itemize}[leftmargin=*]
  \item \textbf{Manivela}: una fila de motor en la unión de revolución (ficha 4) que sigue la velocidad que pide el
        dedo, con impulso limitado a $\tau_{\text{mano}}h$. Si el mecanismo se traba, la manivela no gira.
  \item \textbf{Arrastrar}: un resorte con amortiguamiento crítico entre el dedo y el punto tomado,
        $c = 2\sqrt{km}$, con fuerza máxima.
\end{itemize}

\begin{codigo}[Qué programar]
\begin{tabularx}{\linewidth}{@{}l Y@{}}
\texttt{kit/ruedas.js} & Detectar contacto de discos; fila de rodadura (paralela y transversal); deslizamiento
  regularizado; pérdida por pivoteo.\\
\texttt{kit/bandas.js} & Camino de la banda entre poleas; longitud y ángulos abrazados; tensión neo-Hookeana;
  acople elástico con cabrestante; banda libre entre dos puntos.\\
\texttt{app/} & Herramienta banda; manivela; instrumentos de relación de transmisión en vivo.\\
\end{tabularx}
\end{codigo}

\begin{control}
\begin{enumerate}[leftmargin=*]
  \item ¿Por qué la relación de la transmisión transversal depende de $r$?
  \item ¿Qué se pierde por pivoteo y por qué una rueda delgada pierde menos?
  \item ¿Por qué la elasticidad del caucho es entrópica?
  \item ¿De dónde sale la exponencial del cabrestante?
  \item ¿Por qué una banda floja patina, y qué dos maneras hay de evitarlo?
  \item ¿Por qué la manivela es un motor con par limitado y no una velocidad impuesta?
\end{enumerate}
\end{control}

\begin{cierre}
Ejercicios 6.1--6.5 · pruebas P5 (dos ruedas, $-1/3$), P6 (transversal, $r/R$), P7 (la banda patina en
$\tau_{\max}$) y la longitud de banda · cara pública: ruedas, bandas y manivela.
\end{cierre}

\end{document}
